\scanpage{177}
% archive sheet number: 176
d'une représentation \deleted{dans} appartenant à $\unsure{\mathscr G|\mathcal G}$. \unsure{Alors|Mais},
en vertu de $u*f_{x,x'}=f_{ux,x'}$, \deleted{\unsure{la}}
la représentation de $\Omega$ dans $\Omega*f$ est quotient
de la représentation donnée dans $V$, donc est
dans $\mathscr G$. CQFD.

\figure{\begin{tikzpicture}[x=1cm,y=1cm]
\draw[line cap=round] (0.10,0.19) .. controls (-0.09,0.31) and (-0.06,0.12) .. (0.14,0.10) .. controls (0.37,0.04) and (0.96,0.02) .. (1.02,0.12) .. controls (1.02,0.30) and (0.50,0.42) .. (0.57,0.13) .. controls (0.61,-0.13) and (0.77,-0.25) .. (0.70,-0.40);
\end{tikzpicture}}

Repassons au groupe $G$. \deleted{\unsure{muni d'une}} Il
revient au \deleted{$=$ \unsure{de se donner}} \inserted{\unsure{pour|même}} un sous-espace $A_0$
du \deleted{$\unsure{G^k}$} dual $G^k$ de $G^{(k)}=k(G)$, de
dire que c'est l'ensemble des coefficients
d'un ``groupe de représentations'' \inserted{$\underline{\mathscr G}$} de $G$
(au sens ci-dessus), ou que les
fonctions $f(uv)$ \inserted{sur $G\times G$,} pour $f\in A_0$, sont dans
$A_0\otimes A_0$, et la donnée d'un tel $A_0$
est équivalente à \deleted{\unsure{un tel}} d'un $\underline{\mathscr G}$.
\deleted{\unsure{Si on veut que}} Dire que $\underline{\mathscr G}$ est
\uline{stable par produits tensoriels}, signifie
que $A_0$ est \uline{stable sous multiplication},
i.e. est une algèbre, \unsure{commtve|commutve|commutative} ; dire que
$\underline{\mathscr G}$ est stable par \uline{passage à la contragrédiente}
signifie que \uline{$A_0$ est stable sous $\unsure{\nu|v|\vee}$} ;
